\documentclass[aps,preprint]{revtex4}%
\usepackage{amsfonts}%
\usepackage{amsmath}%
\setcounter{MaxMatrixCols}{30}%
\usepackage{amssymb}%
\usepackage{graphicx}
%TCIDATA{OutputFilter=latex2.dll}
%TCIDATA{Version=4.10.0.2347}
%TCIDATA{CSTFile=revtex4.cst}
%TCIDATA{Created=Saturday, February 08, 2003 11:29:37}
%TCIDATA{LastRevised=Saturday, February 08, 2003 11:37:21}
%TCIDATA{<META NAME="GraphicsSave" CONTENT="32">}
%TCIDATA{<META NAME="DocumentShell" CONTENT="Articles\SW\REVTeX 4">}
\newtheorem{theorem}{Theorem}
\newtheorem{acknowledgement}[theorem]{Acknowledgement}
\newtheorem{algorithm}[theorem]{Algorithm}
\newtheorem{axiom}[theorem]{Axiom}
\newtheorem{claim}[theorem]{Claim}
\newtheorem{conclusion}[theorem]{Conclusion}
\newtheorem{condition}[theorem]{Condition}
\newtheorem{conjecture}[theorem]{Conjecture}
\newtheorem{corollary}[theorem]{Corollary}
\newtheorem{criterion}[theorem]{Criterion}
\newtheorem{definition}[theorem]{Definition}
\newtheorem{example}[theorem]{Example}
\newtheorem{exercise}[theorem]{Exercise}
\newtheorem{lemma}[theorem]{Lemma}
\newtheorem{notation}[theorem]{Notation}
\newtheorem{problem}[theorem]{Problem}
\newtheorem{proposition}[theorem]{Proposition}
\newtheorem{remark}[theorem]{Remark}
\newtheorem{solution}[theorem]{Solution}
\newtheorem{summary}[theorem]{Summary}
\newenvironment{proof}[1][Proof]{\noindent\textbf{#1.} }{\ \rule{0.5em}{0.5em}}
\begin{document}
\preprint{HEP/123-qed}
\title[Short title for running header]{Long title that appears on the title page}
\author{First Author}
\affiliation{My Institution}
\author{Second Author}
\affiliation{My Institution}
\author{Third Author}
\affiliation{Other Institution}
\keywords{one two three}
\pacs{PACS number}

\begin{abstract}
Shell document for REV\TeX{} 4.

\end{abstract}
\volumeyear{year}
\volumenumber{number}
\issuenumber{number}
\eid{identifier}
\date[Date text]{date}
\received[Received text]{date}

\revised[Revised text]{date}

\accepted[Accepted text]{date}

\published[Published text]{date}

\startpage{101}
\endpage{102}
\maketitle
\tableofcontents


\subsection{Ensemble State Generation}

From the preceding we have that $L_{\infty}\left(  \mathsf{Z}^{N}\right)  $
acting on an $a.e.$ nonzero amplitude produces all (unnormalized) pure states,
the question thus arrises on whether $L_{\infty}\left(  \mathsf{Z}^{N}\right)
$ can produce \emph{all} unnormalized states by acting on an a $a.e.$ nonzero
positive kernel $D_{ref}\left(  \mathsf{z}^{N},\mathsf{z}^{\prime N}\right)  $
of a density operator $\in\mathcal{B}_{1+}\left(  \mathcal{H}^{N}\right)  $ in
the following manner%
\begin{equation}
D\left(  \mathsf{z}^{N},\mathsf{z}^{\prime N}\right)  =T\left(  \mathsf{z}%
^{N}\right)  D_{ref}\left(  \mathsf{z}^{N},\mathsf{z}^{\prime N}\right)
T\left(  \mathsf{z}^{\prime N}\right)  ^{\ast}%
\end{equation}
One can start by examining the diagonal elements to obtain
\begin{align}
D\left(  \mathsf{z}^{N},\mathsf{z}^{N}\right)   &  =\left\vert T\left(
\mathsf{z}^{N}\right)  \right\vert ^{2}D_{ref}\left(  \mathsf{z}%
^{N},\mathsf{z}^{N}\right) \nonumber\\
&  \Rightarrow\left\vert T\left(  \mathsf{z}^{N}\right)  \right\vert =\left\{
%
\begin{array}
[c]{c}%
\left(  \frac{D\left(  \mathsf{z}^{N},\mathsf{z}^{N}\right)  }{D_{ref}\left(
\mathsf{z}^{N},\mathsf{z}^{N}\right)  }\right)  ^{\frac{1}{2}};\quad\text{when
}D_{ref}\left(  \mathsf{z}^{N},\mathsf{z}^{N}\right)  \neq0\\
\qquad\qquad0;\qquad\qquad\text{when }D_{ref}\left(  \mathsf{z}^{N}%
,\mathsf{z}^{N}\right)  =0\text{ }%
\end{array}
\right.  \text{ }%
\end{align}
which is well defined $a.e.$ The off diagonal elements will "determine" the
phase of the transformation $T\left(  \mathsf{z}^{N}\right)  =\left\vert
T\left(  \mathsf{z}^{N}\right)  \right\vert e^{i\theta\left(  \mathsf{z}%
^{N}\right)  }$ by
\begin{align}
D\left(  \mathsf{z}^{N},\mathsf{z}^{\prime N}\right)   &  =\left\vert T\left(
\mathsf{z}^{N}\right)  T\left(  \mathsf{z}^{\prime N}\right)  \right\vert
e^{i\left(  \theta\left(  \mathsf{z}^{N}\right)  -\theta\left(  \mathsf{z}%
^{\prime N}\right)  \right)  }D_{ref}\left(  \mathsf{z}^{N},\mathsf{z}^{\prime
N}\right) \nonumber\\
&  \Rightarrow e^{i\left(  \theta\left(  \mathsf{z}^{N}\right)  -\theta\left(
\mathsf{z}^{\prime N}\right)  \right)  }=\frac{D\left(  \mathsf{z}%
^{N},\mathsf{z}^{\prime N}\right)  }{D_{ref}\left(  \mathsf{z}^{N}%
,\mathsf{z}^{\prime N}\right)  \left\vert T\left(  \mathsf{z}^{N}\right)
T\left(  \mathsf{z}^{\prime N}\right)  \right\vert }%
\end{align}
which is well defined as long as
\begin{equation}
\left\vert T\left(  \mathsf{z}^{N}\right)  T\left(  \mathsf{z}^{\prime
N}\right)  \right\vert =0\Leftrightarrow D\left(  \mathsf{z}^{N}%
,\mathsf{z}^{\prime N}\right)  =0
\end{equation}
which means that
\begin{equation}
D\left(  \mathsf{z}^{\prime N},\mathsf{z}^{\prime N}\right)  =0\text{ or
}D\left(  \mathsf{z}^{N},\mathsf{z}^{N}\right)  =0\Rightarrow D\left(
\mathsf{z}^{N},\mathsf{z}^{\prime N}\right)  =0
\end{equation}
when $\breve{D}\geq0$. First we see if is true for proper vectors $\left\vert
\mathsf{v}^{N}\right\rangle =\left\vert \mathsf{v}_{1}\cdots\mathsf{v}%
_{N}\right\rangle $, in that case
\begin{align}
D\left(  \mathsf{v}^{N},\mathsf{v}^{N}\right)   &  =\left\langle
\mathsf{v}^{N}|\breve{D}\mathsf{v}^{N}\right\rangle =\left\langle
\mathsf{v}^{N}|\breve{Q}^{\dagger}\breve{Q}\mathsf{v}^{N}\right\rangle
=0\nonumber\\
&  \Rightarrow\breve{Q}\left\vert \mathsf{v}^{N}\right\rangle =0 \label{c1}%
\end{align}
and
\begin{align}
D\left(  \mathsf{v}^{\prime N},\mathsf{v}^{\prime N}\right)   &  =\left\langle
\mathsf{v}^{\prime N}|\breve{D}\mathsf{v}^{\prime N}\right\rangle
=\left\langle \mathsf{v}^{\prime N}|\breve{Q}^{\dagger}\breve{Q}%
\mathsf{v}^{\prime N}\right\rangle =0\nonumber\\
&  \Rightarrow\breve{Q}\left\vert \mathsf{v}^{\prime N}\right\rangle =0
\label{c2}%
\end{align}
so either condition eq. $\left(  \ref{c1}\right)  $ or eq. $\left(
\ref{c2}\right)  $ clearly implies that
\begin{equation}
\left\langle \mathsf{v}^{N}|\breve{Q}^{\dagger}\breve{Q}\mathsf{v}^{\prime
N}\right\rangle =D\left(  \mathsf{v}^{N},\mathsf{v}^{\prime N}\right)  =0
\end{equation}
To see if this true for the generalized basis we must see if
\begin{equation}
\left\langle \mathsf{z}^{N}|\breve{Q}^{\dagger}\breve{Q}\mathsf{z}%
^{N}\right\rangle =0\Rightarrow\breve{Q}\left\vert \mathsf{z}^{N}\right\rangle
=0
\end{equation}
is true in a rigged Hilbert space \cite{Bohm78}, which we show to be true in
Appendix \ref{rigged}. Thus we have that
\begin{equation}
\left\vert T\left(  \mathsf{z}^{N}\right)  T\left(  \mathsf{z}^{\prime
N}\right)  \right\vert =0\Leftrightarrow D\left(  \mathsf{z}^{N}%
,\mathsf{z}^{\prime N}\right)  =0\text{ if }\breve{D}\geq0
\end{equation}
and we can hence define phase differences as
\begin{equation}
e^{i\left(  \theta\left(  \mathsf{z}^{N}\right)  -\theta\left(  \mathsf{z}%
^{\prime N}\right)  \right)  }=\left\{
\begin{array}
[c]{c}%
\frac{D\left(  \mathsf{z}^{N},\mathsf{z}^{\prime N}\right)  }{D_{ref}\left(
\mathsf{z}^{N},\mathsf{z}^{\prime N}\right)  \left\vert T\left(
\mathsf{z}^{N}\right)  T\left(  \mathsf{z}^{\prime N}\right)  \right\vert
};\left\vert T\left(  \mathsf{z}^{N}\right)  T\left(  \mathsf{z}^{\prime
N}\right)  \right\vert \neq0,D_{ref}\left(  \mathsf{z}^{N},\mathsf{z}^{\prime
N}\right)  \neq0\\
\text{\qquad\qquad}0\qquad\text{\qquad};\text{\qquad otherwise\qquad
\qquad\qquad\qquad\qquad}%
\end{array}
\right.
\end{equation}
which is sufficient to define the transformation $\breve{D}_{ref}%
\rightarrow\breve{D}$ \emph{for any }$\breve{D}$ as long as $D_{ref}\left(
\mathsf{z}^{N},\mathsf{z}^{\prime N}\right)  \neq0\,a.e.$

The transformations $\hat{T}$ can either decrease or leave invariant the rank
of $\breve{D}^{N}$, thus only if $\breve{D}_{ref}$ has maximal rank will all
density operators be generated by such transformations ($\breve{D}_{ref}$ need
not be a density operator associated with a fixed number of particles, all
that is needed is that $\breve{D}_{ref}$ $\in\mathcal{B}_{1+}\left(
\mathcal{H}^{N}\right)  $ with maximal rank and thus could be any density
operator, including ones describing indeterminate number of particles.)

\section{Rigged Hilbert Spaces\label{rigged}}

Rigging a Hilbert space allows one to introduce and use generalized bases in a
mathematical rigorous fashion. Rigged Hilbert spaces are bracketed by a dual
pair of Banach spaces $\mathcal{V}_{L}\supset\mathcal{H\supset}\mathcal{V}%
_{S}$ so that $\mathcal{H}$ and $\mathcal{V}_{S}.$can be naturally embedded in
$\mathcal{V}_{L}.$ The space $\mathcal{V}_{S}$ is dense in $\mathcal{H}$ wrt
to the Hilbert space norm, while both $\mathcal{H}$ and $\mathcal{V}_{S}$ are
dense in $\mathcal{V}_{L}$ wrt the norm of $\mathcal{V}_{L}$. The elements of
a generalized bases $\left\{  \left\vert \mathsf{z}\right\rangle \right\}  $
belong to and generate $\mathcal{V}_{L},$ but they do not belong to
$\mathcal{H}$. The action of an operator $\breve{Q}\in\mathcal{B}\left(
\mathcal{H}\right)  $ on $\mathcal{V}_{L}$ is given by
\begin{equation}
\left(  \breve{Q}\mathsf{z}|\varphi\right)  =\left(  \mathsf{z}|\breve
{Q}^{\dagger}\varphi\right)  \quad\forall\left\vert \varphi\right\rangle
\in\mathcal{V}_{S}%
\end{equation}
where $\breve{Q}^{\dagger}$ is the Hilbert space adjoint of $\breve{Q}$ (which
is well defined as $\breve{Q}$ is a bounded operator), where $\left(
|\right)  $ is the scalar product between $\mathcal{V}_{L}$ and $\mathcal{V}%
_{S}$. Hence one can observe that $\breve{Q}\left\vert \mathsf{z}\right\rangle
=0$ means
\begin{equation}
\left(  \mathsf{z}|\breve{Q}^{\dagger}\varphi\right)  =0\,\forall\left\vert
\varphi\right\rangle \in\mathcal{V}_{S}%
\end{equation}
i.e. $\breve{Q}^{\dagger}\left\vert \varphi\right\rangle =0$ which implies
that $\breve{Q}^{\ddagger}\left\vert \mathsf{z}\right\rangle $ is the zero
vector in $\mathcal{V}_{L}$. Thus $\left\langle \mathsf{z}^{N}|\breve
{Q}^{\dagger}\breve{Q}\mathsf{z}^{N}\right\rangle =0$ or $\left\langle
\mathsf{z}^{\prime N}|\breve{Q}^{\dagger}\breve{Q}\mathsf{z}^{\prime
N}\right\rangle =0\Rightarrow\left\langle \mathsf{z}^{\prime N}|\breve
{Q}^{\dagger}\breve{Q}\mathsf{z}^{N}\right\rangle =0.$


\end{document}